%! Solving Exponential Equations with a Common Base — Unit 7, Lesson 7.3 — Exit Ticket
\documentclass[10pt]{article}
\usepackage{algebra2-article}
\usepackage{algebra2-boxes}
\newcommand{\TallMath}[1]{$\displaystyle #1\rule[-1.4em]{0pt}{3.2em}$}

\begin{document}

\pageheader{Unit 7, Lesson 7.3}{Exit Ticket}
\namedateperiod

\vspace{0.08in}

No notes. On items 1 and 2, the common-base step is part of the answer --- show it.

\begin{enumerate}[leftmargin=1.8em, itemsep=12pt]
  \item \textbf{Solve.} \quad $3^{\,x+1}=81$

        \vspace{4pt}
        Step 1 (one common base): \blank{4.5cm} \quad Step 2 (exponents equal): \blank{3.4cm}

        \vspace{5pt}
        \hfill $x=$ \blank{2.2cm}

  \item \textbf{Solve.} \quad $25^{x}=125$ \quad \emph{(Careful: the two sides do not start with the
        same base.)}

        \vspace{4pt}
        Step 1 (one common base): \blank{4.5cm} \quad Step 2 (exponents equal): \blank{3.4cm}

        \vspace{5pt}
        \hfill $x=$ \blank{2.2cm}

  \item \textbf{Explain, don't compute.} A classmate is trying to solve $2^{x}=-8$ by rewriting $-8$
        as a power of $2$. Explain why no such rewriting exists. Use the \emph{range} of $y=2^{x}$ in
        your answer.

        \vspace{2pt}
        \writelines{2}

  \item \textbf{Multiple choice.} Which equation \emph{cannot} be solved by rewriting both sides with
        a common base?
        \begin{enumerate}[label=\Alph*., leftmargin=1.9em, itemsep=1pt]
          \item $9^{x}=27$
          \item $5^{x}=\frac{1}{125}$
          \item $2^{x}=10$
          \item $8^{x}=\sqrt{2}$
        \end{enumerate}

        \vspace{4pt}
        For the one you chose: does that equation still \emph{have} a solution? Circle one ---
        \textbf{yes} \; / \; \textbf{no} --- and say how you know.

        \vspace{2pt}
        \writeline
\end{enumerate}

\end{document}
